\subsection{Hadits-Hadits Tentang Ilmu}
\begin{itemize}
\item \textit{‘An Anasi bni Mālikin raḍiyallāhu ‘anhu qāla: qāla Rasūlullāhi ṣallallāhu ‘alaihi wa sallam: Talabul-‘ilmi farīḍatun ‘alā kulli muslim.} \\
  \textbf{Artinya: }Dari Anas bin Malik R.A ia berkata: Rasulullah SAW bersabda: ``Menuntut ilmu itu wajib bagi setiap orang muslim.'' \textbf{(HR Ibnu Majah)}

\item \textit{‘An ‘Abdillāhi bni ‘Amri bni al-‘Āṣi raḍiyallāhu ‘anhumā qāla: sami‘tu Rasūlullāhi ṣallallāhu ‘alaihi wa sallam yaqūlu: Innallāha lā yaqbiḍul-‘ilma-ntizā‘an yantazi‘uhū minal-‘ibādi, wa lākin yaqbiḍul-‘ilma biqabḍil-‘ulamā’i, ḥattā iżā lam yubqi ‘āliman-ittakhażan-nāsu ru’ūsan juhhālan, fasu’ilū fa-aftaw bigairi ‘ilmin faḍallū wa aḍallū.} \\
  \textbf{Artinya: }Dari Abdullah bin 'Amr bin al-'Ash R.A ia berkata: Aku mendengar Rasulullah SAW bersabda: ``Sesungguhnya Allah tidak mencabut ilmu dengan menlenyapkan dari dada manusia, tetapi dengan mewafatkan para ulama sehingga setelah tidak ada seorang pun ulama mereka manusia mengangkat orang-orang bodoh menjadi pemimpin. Mereka ditanya, tetapi mereka (pemimpin yang bodoh itu) memberikan petunjuk tanpa ilmu, kemudian tersesatlah mereka dan menyesatkan orang lain pula.'' \textbf{(HR Muslim \& Bukhari)}
\end{itemize}

\subsubsection{Kandungan Isi Ar-Rahman Ayat 33}
\begin{enumerate}
\item Kecanggihan ilmu pengetahuan dan teknologi (IPTEK) harus semakin menumbuhkan kesadaran keimanan kepada Allah SWT itu artinya bahwa semakin luas dan dalamnya ilmu yang dimiliki harus semakin mendekatkan diri kepada-nya, bahwa semuanya merupakan nikmat yang akan dipertanggungjawabkan.
\item Didahulukan penyebutan jin kemudian menusia karena jin lebih memiliki kemampuan menenbus luar angkasa begitu juga perannya di bumi meski lebih terbagtas \textit{(Jin/72 : 29)}.
\item Sebagian ulama menjadikan ayat ini sebagai syarat ilmiah bahwa kekuatan dan penguasaan ilmu menjadi hal yang mutlah dimiliki.
\item Harus disadari bahwa majunya sebuah negara disebabkan besarnya investasi pada kualitas sumber daya manusia (SDM).
\item Allah SWT mengancam kepada jin dan manusia, bahwa kelak di akhirat mereka tidak bisa mengelah akan pertanggungjawaban dari semua nikmat yang sudah diberikan.
\end{enumerate}
